\documentclass[14pt,a4paper]{extarticle}
\usepackage[T2A]{fontenc}
\usepackage[utf8]{inputenc}
\usepackage[english,russian]{babel}
\usepackage{lmodern}
\usepackage{geometry}
\usepackage{array}
\usepackage{tabularx}
\usepackage{longtable}
\usepackage{ltablex}
\usepackage{booktabs}
\usepackage{hyperref}
\usepackage{indentfirst}
\usepackage{microtype}
\usepackage{graphicx}
\usepackage{setspace}
\usepackage{listings}
\usepackage{graphicx}
\usepackage{caption}
\usepackage{float}
\usepackage{tocloft}
\usepackage{titlesec}
\usepackage{newtxtext,newtxmath}

\titleformat{\section}        % форматируем \section
  {\normalfont\bfseries}      % стиль: обычный + жирный
  {\thesection.}              % номер секции с точкой
  {2mm}                       % расстояние между номером и текстом
  {}                           % перед заголовком можно добавить что-то
\titlespacing*{\section}{1.5cm}{*3}{*1}  % {отступ слева}{отступ сверху}{отступ снизу}

\titleformat{\subsection}        % форматируем \section
  {\normalfont\bfseries}      % стиль: обычный + жирный
  {\thesubsection.}              % номер секции с точкой
  {2mm}                       % расстояние между номером и текстом
  {}                           % перед заголовком можно добавить что-то
\titlespacing*{\subsection}{1.5cm}{*3}{*1}  % {отступ слева}{отступ сверху}{отступ снизу}

\captionsetup{justification=centering}
\renewcommand{\ttdefault}{cmtt}
\keepXColumns
\setlength{\parindent}{3em}
\newgeometry{left=2.5cm,right=2.5cm,top=3cm,bottom=2.5cm}

\renewcommand{\contentsname}{\hspace*{4em}{\smallСодержание}}
\setlength{\cftbeforesecskip}{0.5em}
\renewcommand{\cftsecleader}{\cftdotfill{\cftdotsep}}
\renewcommand{\cftsecaftersnum}{.\hspace{2mm}}

\title{Отчёт по "Управление ИТ-проектами"}
\author{Выполнил студент группы М-РИТ-25-1, Станиславчук Сергей}

\begin{document}

\setstretch{1.5}
\noindent
\begin{minipage}[t]{0.01\textwidth}
    \includegraphics{~/gitrepos/LSTU_MAG/5_2/LSTU.jpg}
\end{minipage}%
\begin{minipage}[c]{1.1\linewidth}
    \centering
    {\small
    МИНИСТЕРСТВО НАУКИ И ВЫСШЕГО ОБРАЗОВАНИЯ РФ \\[1em]
    ФЕДЕРАЛЬНОЕ ГОСУДАРСТВЕННОЕ БЮДЖЕТНОЕ \\
    ОБРАЗОВАТЕЛЬНОЕ УЧРЕЖДЕНИЕ ВЫСШЕГО ОБРАЗОВАНИЯ \\
    «ЛИПЕЦКИЙ ГОСУДАРСТВЕННЫЙ ТЕХНИЧЕСКИЙ УНИВЕРСИТЕТ»
    }
\end{minipage}%

\vspace{3em}

\noindent
\hspace{1em}%
\begin{minipage}[t]{0.8\textwidth}
\textbf{Институт} \hspace{2em} компьютерных наук \\
\textbf{Кафедра} \hspace{2em} автоматизированных систем управления
\end{minipage}


\vspace{6em}

\begin{center}
\textbf{Практическая работа №13} \\
по дисциплине \\ <<Управление ИТ-проектами>>
\end{center}

\vspace{6em}

\noindent
\begin{tabular*}{\textwidth}{@{\extracolsep{\fill}} l c r }
    Студент М-РИТ-25-1 & \makebox[4cm]{\hrulefill} & Станиславчук С.М. \\[2em]
Руководитель & \\
К.т.н., доцент & \makebox[4cm]{\hrulefill} & Ведищев В.В. \\
\end{tabular*}

\vfill
\begin{center}
    Липецк, 2026
\end{center}

\thispagestyle{empty}

\newpage
\tableofcontents
\newpage

\section{Формирование и поддержка конфигурации программного продукта}

\subsection{Понятие «управление конфигурацией»}

Управление конфигурацией представляет собой совокупность процессов и методов, направленных на контроль изменений программного продукта на протяжении всего жизненного цикла. В рамках разработки системы корпоративного обмена сообщениями управление конфигурацией обеспечивает согласованность исходного кода, документации, зависимостей и сборочных артефактов.

Конфигурация включает в себя все компоненты системы: исходный код серверной и клиентской частей, файлы сборки (CMake), настройки окружения, а также документацию и тесты. Основной задачей управления конфигурацией является обеспечение воспроизводимости сборки и контроля всех изменений.

\subsection{Процесс управления конфигурацией}

Процесс управления конфигурацией включает идентификацию конфигурационных единиц, контроль изменений, учёт версий и аудит состояния системы. На этапе идентификации определяется перечень элементов, подлежащих контролю, включая исходные файлы, зависимости и настройки.

Контроль изменений осуществляется через систему контроля версий, где каждое изменение фиксируется в виде коммита. Все изменения проходят проверку перед включением в основную ветку проекта. Учёт версий реализуется через систему тегирования и ветвления, что позволяет фиксировать стабильные состояния системы.

Аудит конфигурации проводится для проверки соответствия текущего состояния системы установленным требованиям и документации.

\subsection{Планирование и настройка управления конфигурацией}

Планирование управления конфигурацией включает определение структуры репозитория, правил ветвления и процедур внесения изменений. В рамках проекта используется модель с основными ветками \texttt{main} и \texttt{develop}, а также вспомогательными ветками \texttt{feature/*}, \texttt{bugfix/*} и \texttt{release/*}.

Настройка включает определение правил именования веток, формата сообщений коммитов, а также процедур code review. Для сборки проекта используется CMake, обеспечивающий единообразие конфигурации сборки на разных платформах.

Дополнительно настраиваются CI/CD-процессы, обеспечивающие автоматическую сборку и тестирование системы при внесении изменений.

\subsection{Выполнение контроля конфигурации}

Контроль конфигурации заключается в управлении изменениями и предотвращении неконтролируемого внесения правок. Все изменения вносятся через отдельные ветки и проходят процедуру проверки перед объединением с основной веткой.

Использование системы контроля версий позволяет отслеживать историю изменений, выявлять источники ошибок и при необходимости выполнять откат к предыдущим версиям. Важным элементом контроля является привязка изменений к конкретным задачам (тикетам), что обеспечивает прослеживаемость.

\subsection{Мониторинг состояния конфигурации}

Мониторинг состояния конфигурации направлен на отслеживание текущего состояния системы и выявление отклонений. В рамках проекта используется контроль статуса сборок, результатов тестирования и состояния веток в репозитории.

Дополнительно осуществляется анализ логов сборки и выполнения системы, что позволяет выявлять ошибки на ранних этапах. Регулярный мониторинг обеспечивает своевременное обнаружение проблем и поддержание стабильности программного продукта.

Таким образом, управление конфигурацией обеспечивает контроль изменений, воспроизводимость сборки и устойчивость разработки в условиях командной работы.

\end{document}
